\documentclass[11pt]{amsart}
\usepackage[margin=1in]{geometry}
\usepackage{amsmath,amssymb,mathtools}
\usepackage[colorlinks=true,linkcolor=blue,citecolor=blue,urlcolor=blue]{hyperref}

\newcommand{\E}{\mathbf E}
\newcommand{\R}{\mathbb R}
\newcommand{\op}{\mathrm{op}}

\title{The Gaussian Row--Column Equivalence in arXiv:1203.3713\\
Follows by Symmetric Block Dilation}
\author{agent\_lane\_03}
\date{27 August 2026}

\begin{document}
\maketitle

\section*{Status}

\textbf{Status: literature\_implied\_answer (full Gaussian question).}
The implication below completely proves the Gaussian equivalence suggested
by Riemer--Sch\"utt.  It is not claimed as a new result: it is a direct
reformulation of a later theorem of Lata\l a--van Handel--Youssef.

\section*{Source question}

Let $A=(a_{ij})\in\R^{n\times n}$ be deterministic, let $(g_{ij})$ be
independent standard Gaussian variables, and put $G=(a_{ij}g_{ij})$.  Define
\[
 R(G)=\max_i\Big(\sum_jG_{ij}^2\Big)^{1/2},\qquad
 C(G)=\max_j\Big(\sum_iG_{ij}^2\Big)^{1/2}.
\]
On page 3 of arXiv:1203.3713, immediately after display (2), Riemer and
Sch\"utt write that they know no matrix for which $\E\|G\|_{\op}$ is not of
the same order as that display, namely
\begin{equation}\label{eq:question}
                    \E R(G)+\E C(G).
\end{equation}
Thus the suggested dimension-free equivalence is
\[
                 \E\|G\|_{\op}\asymp \E R(G)+\E C(G),
\]
with universal constants, uniformly over $n$ and $A$.

\section*{Supporting theorem}

Lata\l a--van Handel--Youssef prove the following in
\cite[Theorem 3.1, p.~16]{LVHY}; it is the probabilistic form of their main
Theorem 1.1.  If $X$ is a symmetric matrix whose upper-triangular entries
are independent centered Gaussian variables, then, for every
$2\le p\le\infty$,
\[
 \E\|X\|_{S_p}\asymp
 \E\left[\left(\sum_i\left(\sum_jX_{ij}^2\right)^{p/2}\right)^{1/p}\right].
\]
At $p=\infty$ this reads
\begin{equation}\label{eq:LVHY}
        \E\|X\|_{\op}\asymp
        \E\max_i\Big(\sum_jX_{ij}^2\Big)^{1/2},
\end{equation}
with universal constants.

\section*{Identification and complete deduction}

Form the symmetric $2n\times2n$ block dilation
\[
                 X=\begin{pmatrix}0&G\\G^{\mathsf T}&0\end{pmatrix}.
\]
The entries in the upper-right block are precisely the independent Gaussian
variables $a_{ij}g_{ij}$; all other upper-triangular entries are deterministic
zeros.  Hence $X$ satisfies the hypotheses of \eqref{eq:LVHY}.  Moreover,
\[
 X^2=\begin{pmatrix}GG^{\mathsf T}&0\\0&G^{\mathsf T}G\end{pmatrix},
 \qquad \|X\|_{\op}=\|G\|_{\op},
\]
and inspection of its two row blocks gives
\[
       \max_{k\le2n}\Big(\sum_{\ell\le2n}X_{k\ell}^2\Big)^{1/2}
       =\max\{R(G),C(G)\}.
\]
Applying \eqref{eq:LVHY} therefore yields
\begin{equation}\label{eq:max}
              \E\|G\|_{\op}\asymp \E\max\{R(G),C(G)\}.
\end{equation}
Finally, pointwise,
\[
 \max\{R,C\}\le R+C\le2\max\{R,C\}.
\]
Taking expectations in this inequality and combining with \eqref{eq:max}
proves
\[
 \boxed{\displaystyle
 \E\|G\|_{\op}\asymp \E R(G)+\E C(G)}.
\]
This is exactly the equivalence suggested after display (2) of
arXiv:1203.3713.

\section*{Provenance, scope, and search evidence}

The supporting authors explicitly state that their $p=\infty$ theorem
settles Lata\l a's symmetric Gaussian conjecture.  In the inspected source,
they do not explicitly identify the Riemer--Sch\"utt row--column formulation;
the rectangular statement is obtained here by the standard block-dilation
identification.  The correct status is therefore
\texttt{literature\_implied\_answer}, not a new proof of the run.

The result covers every deterministic coefficient matrix with independent
standard Gaussian multipliers.  It does not settle the source paper's
separate estimate for arbitrary independent mean-zero entries.  The bounded
status search covered the run indexes for arXiv:1203.3713 and the phrases
``Gaussian matrix'', ``row norm'', ``column norm'', and ``Lata\l a'', then
checked the exact arXiv sources 1203.3713 and 1711.00807.  A prior run packet
for Lata\l a's symmetric conjecture led to the decisive Theorem 3.1.

\begin{thebibliography}{9}
\bibitem{RS}
S.~Riemer and C.~Sch\"utt,
\emph{On the Expectation of the Norm of Random Matrices with
Non-Identically Distributed Entries}, arXiv:1203.3713 (2012),
\url{https://arxiv.org/abs/1203.3713}.

\bibitem{LVHY}
R.~Lata\l a, R.~van Handel, and P.~Youssef,
\emph{The dimension-free structure of nonhomogeneous random matrices},
Invent. Math. \textbf{214} (2018), 1031--1080; arXiv:1711.00807,
\url{https://arxiv.org/abs/1711.00807}.
\end{thebibliography}

\end{document}
